% 原书分章PDF第6—8页：第30.3节全部9段，长引文完整。
\section{遭到侵蚀的模型、冯·诺依曼与理论的终结}

% S30.3.P01
匈牙利物理学家约翰·冯·诺依曼不仅是计算机科学的先驱、博弈论的奠基者，而且与爱因斯坦（Einstein）、海森堡（Heisenberg）等同行一样，积极关注哲学与科学的关系。他对认识论的关切，可以从其文章“物理科学中的方法”（Method in the Physical Sciences）中得到印证；该文发表于1955年，并于2000年收入《诺依曼文集》（The Neumann Compendium）重印。在一段经常被引用的文字中，这位物理学家几乎没有给误解其立场留下余地：“科学并不试图解释，甚至几乎不试图阐释，它们主要是在构造模型。所谓模型，是一种数学构造；在附加某些言语阐释之后，它描述所观察到的现象。这样一种数学构造的证成，唯一而且恰恰在于，人们预期它能发挥作用——也就是，能够正确地描述相当广泛范围内的现象。”

% S30.3.P02
对于某些科学家来说，这段简短文字或许不会引起多少不安或困惑（Starmans，2011a）。其中呈现的观点，似乎至多与日常实践中的经验部分相容。任何物理学家、生物学家、社会科学家、经济学家、气候科学家或工程师，都会认同模型在科学中无处不在这一主张。此外，模型也明显存在于科学之外的专业实践和日常生活之中。然而，从认识论的角度看，冯·诺依曼的这段话在多个方面都存在问题。

% S30.3.P03
首先，他主张把科学几乎完全还原为模型构造。模型并非只是在理论化过程中提供解释框架，或附加的启发价值，而是整个科学事业的本质；两者似乎或多或少是等价的。这个命题在多大程度上构成问题，取决于所用模型概念的地位与范围。在这一点上，冯·诺依曼同样十分明确。首先，他主要把科学模型视为附有言语说明的数学构造。在今天，这一说法似乎很容易受到反驳，因为科学技术使用的不仅是符号化的数学模型，还包括大量有形的、物理的、类比的模型，以及图形模型或数字仿真模型。而且，冯·诺依曼完全以功利用途来证成模型的使用；说明某种（潜在的）实在，似乎不再是目标，追求真理也不再是目标。这既体现了极端的实用主义，也体现了一种几乎毫不掩饰、且必然随之而来的工具主义（instrumentalism）与反实在论（antirealism）。它还揭示了一种科学观：解释——无论是因果解释还是其他解释——不再是主要的研究功能，真理已经成为一个孤立、甚至无用的概念。这与前文引用过的统计学家乔治·博克斯的观点相近。

% S30.3.P04
更重要的是，大数据与数据科学近来的发展，使上述思考可以被看作一种前奏，或者使其更有理由（a fortiori）获得重要性；这些发展对科学哲学和方法论具有更加深远的后果。例如，可以想到科学记者克里斯·安德森（Chris Anderson）一篇声名狼藉却经常被引用的文章；他早在2008年就宣告了科学方法论与科学知识的终结。他的“理论的终结：数据洪流使科学方法过时”（The End of Theory: The Data Deluge Makes the Scientific Method Obsolete）（Anderson，2008），对于哲学家、方法论研究者和统计学家同样重要。略选几段安德森有时相当平实的评论，便足以说明问题。

% S30.3.P05
首先，这位作者声称，谷歌等公司实际上正在筛检文明诞生以来测量最充分的时代，“把这个庞大的资料总体视为人类处境的实验室”，建立一个由“拍字节时代的孩子们”（children of the Petabyte Age）组成的社会。显然，作者并不回避任何煽情的表达。他不是未来学家，而是毫不含糊地声称，所描绘的情形与其说是遥远的前景，不如说就是现状。按照安德森的说法，在数据科学时代，信息处理不再涉及简单的低维分类、自然规律和理论，而涉及他所谓“不问维数的统计学”（dimensionally agnostic statistics）。这要求一种本质上不同的方法论，最终充分尊重数据，使它成为可以“作为整体加以可视化”的东西。为此，我们必须“首先从数学角度看待数据，然后再为其建立语境”。

% S30.3.P06
这位作者表现得像新学说的一位真正辩护士。按照安德森的说法，谷歌的成功建立在这样一种信念之上：我们并不理解为什么这个网页比另一个网页更好，而且实际上也不必为此费心；完全依赖统计就够了，“不需要任何语义分析或因果分析”。在这位作者看来，这几乎涵盖了人类行动的任何方面或维度：“把从语言学到社会学的一切人类行为理论统统抛开。忘掉分类学、本体论和心理学。谁知道人们为什么做他们做的事？关键在于，他们做了，而我们能够以前所未有的保真程度追踪和测量这些行为。只要有足够的数据，数字就会自己说话。”确实，我们把科学理解为提出假设、建立模型与进行检验的整套观念，似乎正在迅速过时。

% S30.3.P07
安德森断言，物理学领域似乎就是一个很好的例子。首先，牛顿定律被证明是真理的粗糙近似；它们有用，但在原子层面是错的。随后，在20世纪初，量子力学又提供了另一个模型；这个模型诚然有所改进，却仍有缺陷，仍然是对现象背后更复杂实在的误表征或歪曲。此后，生物学也受到审判。高中教科书中基于孟德尔遗传学的“显性”与“隐性”基因的经典叙述，似乎比经典力学更加粗糙地过度简化了实在。事实上，“基因—蛋白质相互作用及表观遗传学其他方面的发现，已经挑战了把 DNA 视为命运的观点，甚至提供了环境能够影响可遗传性状的证据；而这曾被认为在遗传学上是不可能的”。更糟的是，关于生物现象的真实知识越多，提出能够解释这些现象的模型的可能性，反而必然越小。拍字节时代的孩子们将以一种全然不同的方式感知和理解世界。科学进步不再需要任何机制性解释，因果模型将被相关模式取代，大统一理论也不再被需要或成为追求的目标。

% S30.3.P08
一年后，谷歌研究者 Alon Havely、Peter Norvig 和 Fernando Pereira 发表了“数据不可思议的有效性”（The Unreasonable Effectiveness of Data）。这些作者表示，当研究对象是人而不是基本粒子时，简单的数学公式与优雅的理论作用有限；参数统计模型也同样如此。要理解和预测人的行为，更可以依靠许多拍字节的粗糙、未经编目、非结构化，必要时甚至失真的数据。这些数据先用高效的数据结构／架构表示出来，再由具有模式识别能力的智能算法锻造成知识。在 Starmans（2016b，2017b）中，我们曾试图更公允地呈现上述作者的观点；不过，对其思想的这一总体描述，已经足以说明统计学与“科学方法”目前的处境，以及它们背后的哲学如何显得承受着压力。

% S30.3.P09
此外，安德森与 Havely 等人的观点，在大数据与数据科学的大众及专业文献中几乎随处可见，尽管有时会以稍为缓和的形式出现。显然，这并不意味着不存在相反的、更加细致的立场。谷歌也不例外（Starmans，2015a），尤其是在当代关于人工智能风险，以及更一般地关于人工智能与数据科学道德维度的争论中。例如，数学家、后来转而怀疑数据科学的凯茜·奥尼尔（Cathy O'Neil）于2016年出版的《数学毁灭武器：大数据如何加剧不平等并威胁民主》（Weapons of Math Destruction; How Big Data Increases Inequality and Threatens Democracy）近来大获成功，就是一个说明。无论如何，本节所描绘的发展，暗示了一种反认识论与反统计学立场，似乎只给统计学和认识论留下了漏洞百出的前景；这或许会促使人们考虑双方联手的问题。然而，这显然远非易事。
